\section{Experimental Setup}
\label{sec:experiments}
\subsection{Historical data and evaluation scope}
We audit fixed scored test cases from three historical runs: ASSISTments interaction correctness ($n=105107$), OULAD student success ($n=6519$) and a synthetic binary stress test ($n=1000$). ASSISTments uses $y=1$ for a correct response and school-ID parity as a synthetic grouping proxy, not a protected attribute. OULAD uses $y=1$ for Pass/Distinction and binary gender for auditing. Appendix~\ref{app:datasets} gives task and split details. These are outcome-prediction testbeds; descending success scores do not prioritize the highest-risk students for intervention.

The historical runner used a seeded shuffled 80/20 split and seed 42. We do not claim a temporal early-warning horizon, student-cluster independence for interaction rows, or prospective treatment effects. The revision replays emitted scores without retraining historical models. Policy search uses labels from the same retrospective audit sample, so selected operating points are exploratory, not an independently validated final policy.

\subsection{Policy domains}
All datasets use the threshold grid $\{.01,\allowbreak .05,\allowbreak .10,\allowbreak .15,\allowbreak .20,\allowbreak .30,\allowbreak .40,\allowbreak .50,\allowbreak .60,\allowbreak .70,\allowbreak .80,\allowbreak .90,\allowbreak .95,\allowbreak .99,\allowbreak .995,\allowbreak .999\}$ bits. Capacity grids are ASSISTments $\{100,\allowbreak 250,\allowbreak 500,\allowbreak 1000,\allowbreak 2000\}$, OULAD $\{50,100,200,500,1000\}$ and synthetic $\{20,50,100,200,300\}$. Each full grid has 80 candidates. Fixed-budget recovery uses only the 16 thresholds at $K=1000,500,100$, respectively. Full-grid profile comparisons are reported separately. We check gap $\leq.05$ and automation coverage $\geq.50$ and apply the algorithm in Section~\ref{sec:method}.

\subsection{Exploratory contestation study}
\label{sec:human-study}
The original study has 16 volunteers, eight standardized appeal scenarios and 64 case responses per condition. Profiles include engineering/data science, analytics and audit/compliance. Each participant responds to four cases under A (logs/cards) and four under B (structured per-case evidence, policy provenance and replay narratives). The author confirmed that G1 receives odd cases in A and even cases in B, with G2 reversed. Each case therefore has eight responses per condition. No participant sees the same case in both conditions. Detailed recruitment procedures, compensation and training records were not available in the inspected release.

For each case, participants choose a diagnostic label and remediation action, cite evidence, and rate confidence from 1 to 5. Evidence length is counted in Unicode characters, not words or a quality rubric. Time was self-reported for the questionnaire; we do not estimate per-case staff savings. The case materials include hypothetical exceptions and feature changes beyond the fixed-score replay engine. Consequently this study evaluates availability of structured evidence as a bundle, not execution of every scenario by the software or the isolated causal effect of replay.

\paragraph{Consent and ethics reporting.}
The author confirmed voluntary participation and authorization. The forms contain a consent item. No formal ethics-committee approval or exemption was obtained; we report this explicitly rather than infer an institutional authorization. No participant-level responses or free-text evidence are redistributed in the revised public analysis outputs. The setting and documentation limit the strength and generalizability of the study.

\paragraph{Retrospective analysis.}
We recompute participant-level condition averages from the original CSVs and case scoring keys. The revision uses two-sided signed-rank tests by enumerating all sign assignments to nonzero paired midranks; zero differences are excluded and equal absolute differences receive average ranks. No normal approximation or continuity correction is used. Holm adjustment covers diagnosis, remediation, confidence and evidence length. Paired mean differences have percentile 95\% intervals from 10000 participant-bootstrap replicates (seed 42). This analysis was specified after inspecting historical results and is not preregistered. Diagnosis and remediation correctness coincide on all 128 case responses and are not independent replications. Software versions and exact counts are recorded with the released revision scripts.
